% ============================================================
% PAGE 4 — Proposed Method: FusionDet Overview and Backbone
% ============================================================

\section{Proposed Method: FusionDet}
\label{sec:method}

FusionDet is a single-stage, anchor-free object detector
designed to advance the speed--precision Pareto frontier by
combining three architectural innovations: (i) a
\textit{reconfigured lightweight backbone} that maximizes
representational capacity per floating-point operation; (ii) a
\textit{bidirectional attention-augmented neck} that fuses
multi-scale features more effectively than standard FPN or
PANet; and (iii) a \textit{decoupled three-branch head} trained
with CIoU regression loss and focal classification loss.
% TODO: Insert architecture figure; label=fig:architecture
Figure~1 (omitted in this preprint)
illustrates the complete pipeline.

At a high level, the detector processes an input image
$\mathbf{I} \in \mathbb{R}^{H \times W \times 3}$ through the
backbone to extract four feature maps at progressively coarser
spatial resolutions; the neck fuses these maps bidirectionally
and applies channel-and-spatial attention recalibration; and
the head generates, at each spatial location of each fused
feature map, a classification probability vector, four
bounding-box regression offsets, and a scalar objectness score.
The following subsections describe each component in detail.

\subsection{Module 1: Reconfigured Backbone}

\subsubsection{Design Rationale}

The backbone is responsible for transforming a raw image into
a hierarchy of progressively abstract feature representations.
In the context of real-time detection, the ideal backbone
achieves the highest possible semantic discriminability per
multiply-accumulate (MAC) operation, minimizing the inference
budget consumed before the neck and head stages.

Standard ResNet variants~\cite{he2016deep} offer strong
representational capacity but are over-parametrized for
detection: their full width is designed for ImageNet
classification, where the final feature map is pooled to a
single vector, discarding spatial information.  Detection
requires spatially dense features across multiple scales;
consequently, a narrower backbone with well-placed skip
connections can match or exceed ResNet performance at a
fraction of the computational cost.

Lightweight architectures such as
MobileNetV2~\cite{sandler2018mobilenetv2} and
MobileNetV3~\cite{howard2019searching} achieve extreme
efficiency via depthwise-separable convolutions, which
factorize a standard $k \times k$ convolution into a spatial
depthwise step (applied independently per channel) and a
pointwise $1 \times 1$ step (mixing channels).  For a $k
\times k$ convolution from $C_{in}$ to $C_{out}$ channels over
a feature map of spatial size $H' \times W'$, the operation
count is:
\begin{equation}
    \Omega_{\text{std}} =
    k^2 \cdot C_{in} \cdot C_{out} \cdot H' \cdot W'
    \label{eq:std_conv}
\end{equation}
while its depthwise-separable counterpart requires:
\begin{equation}
    \Omega_{\text{dws}} =
    \left(k^2 \cdot C_{in} + C_{in} \cdot C_{out}\right)
    \cdot H' \cdot W'
    \label{eq:dws_conv}
\end{equation}
For $k = 3$ and $C_{in} = C_{out} = C$, the ratio
$\Omega_{\text{dws}} / \Omega_{\text{std}} \approx 1/9 +
1/C \approx 1/9$, delivering roughly an $8\times$ reduction in
MACs with well-managed accuracy loss.

\subsubsection{Backbone Architecture}

FusionDet's backbone is a modified MobileNetV3-Large variant
with three targeted reconfiguration choices:

\textbf{Stem convolution.}  The standard $3\times3$ stride-2
stem is retained, projecting the three-channel input to 16
channels and halving spatial resolution.  An additional
hard-swish activation~\cite{howard2019searching} is applied
before the first inverted residual block to improve gradient
flow in the early layers.

\textbf{Inverted residual blocks.}  The backbone comprises 15
inverted residual blocks arranged in four stages with strides
$\{1, 2, 2, 2\}$, producing output feature maps at strides
$\{8, 16, 32, 64\}$ relative to the input, denoted $P_2, P_3,
P_4, P_5$.  These four stages constitute the pyramid levels
fed into the neck.  Compared with the original MobileNetV3-Large,
the expansion ratios of the final two stages are increased from
6 to 8 to improve feature richness at the coarser, more
semantically meaningful levels.

\textbf{Partial Channel Fusion (PCF) blocks.}  Inspired by the
Cross-Stage Partial (CSP) design~\cite{bochkovskiy2020yolov4},
we insert PCF blocks at the transition points between the third
and fourth backbone stages.  Each PCF block splits the input
feature map channel-wise into two equal halves; one half passes
through two stacked depthwise-separable convolutions while the
other bypasses them, and both halves are concatenated and
projected back to the original channel count.  PCF blocks add
a residual-style shortcut that preserves gradient flow and
increases the effective receptive field of the backbone without
introducing additional parameters proportional to channel
count squared.

\subsubsection{Output Dimensions}

The four pyramid-level feature maps produced by the backbone
and passed to the neck have the following dimensions for a
reference input of $640 \times 640$:
\begin{align}
    P_2 &\in \mathbb{R}^{80 \times 80 \times 40}  \label{eq:p2} \\
    P_3 &\in \mathbb{R}^{40 \times 40 \times 80}  \label{eq:p3} \\
    P_4 &\in \mathbb{R}^{20 \times 20 \times 160} \label{eq:p4} \\
    P_5 &\in \mathbb{R}^{10 \times 10 \times 320} \label{eq:p5}
\end{align}
The channel doubling at each stride halving follows the
standard convention; the absolute channel counts are chosen to
match the capacity of YOLOv8-s while reducing total GFLOPs by
approximately $18\%$ on the backbone stage alone.

\subsection{Backbone Ablation: Efficiency Analysis}

To validate the reconfiguration choices, we compare three
backbone variants against the reference YOLOv8-s backbone:
(a) vanilla MobileNetV3-Large with no modifications, (b)
MobileNetV3-Large with expanded final-stage width only, and (c)
the full reconfigured backbone with expanded width and PCF
blocks.  Table~\ref{tab:backbone_ablation} summarizes
GFLOPs, parameter counts, and the mAP achieved when each
backbone is combined with the same standard FPN neck and
decoupled head.

\begin{table}[!t]
\centering
\caption{Backbone efficiency analysis on MS-COCO val2017.}
\label{tab:backbone_ablation}
\renewcommand{\arraystretch}{1.15}
\setlength{\tabcolsep}{4pt}
\begin{tabular}{l c c c}
\hline\hline
\textbf{Backbone}          & \textbf{GFLOPs} &
\textbf{Params (M)} & \textbf{mAP} \\
\hline
MobileNetV3-L (vanilla)    & 8.1  & 5.4  & 41.3 \\
MobileNetV3-L + wide       & 9.7  & 6.2  & 42.8 \\
\textbf{FusionDet backbone}& \textbf{10.3} & \textbf{6.8} &
\textbf{44.1} \\
YOLOv8-s backbone (ref.)   & 12.6 & 8.3  & 44.0 \\
\hline\hline
\end{tabular}
\end{table}

The full reconfigured backbone achieves accuracy on par with
the larger YOLOv8-s backbone while consuming $18.3\%$ fewer
GFLOPs, confirming that PCF blocks and expanded-stage width
together recover the accuracy gap of vanilla MobileNetV3-Large
at low computational overhead.
